\documentclass[10pt,a4paper]{amsart}
\usepackage[margin=19mm]{geometry}
\usepackage{amsmath,amssymb,mathtools}
\usepackage[scr=rsfs,cal=euler]{mathalfa}
\usepackage{xfrac}
\usepackage[unicode,hidelinks,bookmarks=false]{hyperref}
\newcommand{\ie}{i.e.\ }
\newcommand{\cf}{cf.\ }
\newcommand{\resp}{respectively\ }
\newcommand{\Subsection}{\subsection}
\providecommand{\nobreakdash}{\nobreak\hbox{-}}
\setlength{\emergencystretch}{2em}
\multlinegap0pt
\usepackage{bm}
\usepackage{bbm}
\usepackage{dsfont}
\usepackage{xspace}
\usepackage{cancel}
\usepackage{mathtools}
\usepackage{amsthm}
\makeatletter
\@ifundefined{theorem}{\newtheorem{theorem}{Theorem}}{}
\makeatother
\newcommand{\Soc}{\operatorname{Soc}} 
\title[Decompositions into a direct sum of projective and stable su]{Decompositions into a direct sum of projective and stable submodules}
\author{Gulizar Gunay, Engin Mermut}
\date{Source: arXiv:2503.07271v3 (CC BY 4.0)}
\begin{document}
\maketitle

\noindent\emph{Source and licence.} Decompositions into a direct sum of projective and stable submodules. Version arXiv:2503.07271v3 (CC BY 4.0), distributed under the Creative Commons Attribution 4.0 International licence, as recorded in the arXiv licence metadata for this exact version and shown on the abstract page https://arxiv.org/abs/2503.07271v3. Full rights notice: licenses/arxiv-2503.07271-CC-BY-4.0.txt.\par\smallskip

\noindent\textit{Editorial scope.} Counted unit: the Theorem whose source label is \texttt{Ex:Cyclicmodule} in the article (source0.tex lines 575--576, source label \texttt{Ex:Cyclicmodule}), delivered together with the article's own complete original proof of it (source0.tex lines 577--591, one \texttt{proof} environment of 15 lines). No mathematics is written, repaired or re-derived: statement and proof are the article's own text line for line. Required context retained uncounted: none. Every definition, display and label of those lines is retained unchanged. Omitted from this package and not redistributed: 1--574, 592--786. Nothing inside a retained excerpt is omitted or altered: every retained body is delivered exactly as the article's lines are written, so no omission receipt of any kind accompanies this package. Citations: the retained excerpts carry no citation command at all, so no bibliography entry of the article is transcribed and no reference is redistributed. Reference adaptation: none was needed and none was made; every reference inside the delivered unit resolves inside it, so no unresolved reference or citation is produced. Printed slips kept literal: no printed word, symbol, hypothesis or number of the counted statement or of its proof was altered. Layout and class: the article's own class is \texttt{amsart} with packages including amsmath, amssymb, comment, fullpage, xy, natbib, color, amsthm, hyperref; the wrapper is the lane's pinned \texttt{amsart} editorial wrapper at 10pt on A4, which restates the theorem-like environments the retained text needs and adds the article's own macro definitions that the retained text needs (\texttt{\textbackslash Soc}), with extra wrapper packages (bm, bbm, dsfont, xspace, cancel, mathtools). The delivered numbering is the unit's own. Assets: no retained span contains a figure environment, a graphics-inclusion command, a PDF-page inclusion, a table or a TikZ picture; no image file is copied into this package, the package declares no asset dependency and no image file is redistributed with it. Licence: version arXiv:2503.07271v3 of this article is distributed under the Creative Commons Attribution 4.0 International licence, as recorded in the arXiv licence metadata for this exact version and shown on the abstract page https://arxiv.org/abs/2503.07271v3. A full rights notice is delivered at licenses/arxiv-2503.07271-CC-BY-4.0.txt. Characters: every retained excerpt is pure ASCII, so the delivered engine, XeLaTeX with its default Latin Modern text font, reports no missing character.
\begin{theorem} \label{Ex:Cyclicmodule} If $R$ is a right semiartinian right V-ring that is not a semisimple ring, then there exists a \emph{cyclic} (right) $R$-module $M$ that cannot be decomposed as $M=P\oplus N$ such that $P$ is projective and $N$ is stable.
\end{theorem}

\begin{proof}

Since $R$ is a right semiartinian ring that is not a semisimple ring, the nonzero right $R$-module $R/\Soc(R_R)$ must have a simple submodule $C/\Soc(R_R)$ where $\Soc(R_R)\subseteq C \subseteq R$. Let $c\in C\setminus \Soc(R_R)$. Then \[C/\Soc(R_R)=(c+\Soc(R_R))R=(cR+\Soc(R_R))/\Soc(R_R).\] Let $D=cR$. Note that $D$ is not semisimple; otherwise, $D=cR\subseteq \Soc(R_R)$, contradicting $c \notin \Soc(R_R)$.
Thus $\operatorname{Soc}(D) \neq D$. This implies that $\Soc (D)$ is not finitely generated. Indeed, if it were finitely generated, it would be a direct sum of finitely many simple right
$R$-modules, which are injective as $R$ is a right V-ring, and so $\Soc(D)$ would be injective, which would then be a direct summand of $D$. But this would lead to a contradiction since $\Soc (D)$ is an essential submodule of $D$ (as the ring $R$ is right semiartinian).

Consider the module $D/\Soc(D)$; it is simple since it is isomorphic to the simple module $ C/\Soc(R_R) $:
\[D/\Soc(D)=cR/\Soc(cR)=cR/cR\cap \Soc(R_R)\cong (cR+\Soc(R_R))/\Soc(R_R)= C/\Soc(R_R).\] {Therefore, $\Soc(D)$ is a maximal submodule of $D$. Since the semisimple module} $\Soc(D)$ is not finitely generated, there exists a decomposition $\Soc(D)=A\oplus B$ where both $A$  and $B$ are semisimple submodules of $D$ that are not finitely generated. Let $M=D/A$.

Suppose that $M=D/A=(P/A)\oplus (N/A)$ for some submodules $P$, $N$ of $D$ such that $A\subseteq P$, $N \subseteq D=cR\subseteq R$, where $P/A$ is projective and $N/A$ {is stable}.

Firstly, we must have $ P\neq D $; otherwise, $M=D/A=P/A$ would be projective, making $ A $  a direct summand of the cyclic module $P=D=cR$, and contradicting the fact that $ A $ is not finitely generated. Suppose that $P\supseteq \Soc(D)$. Since $\Soc(D)\subseteq P \subseteq D$ and $\Soc(D)$ is a maximal submodule of $D$, we must have either $P=\Soc(D)$ or $P=D$. As seen above $P\neq D$, and so we must have $P=\Soc(D)=A\oplus B$. In this sum, $B\cong  P/A$ is a cyclic $R$-module since it is a quotient of the cyclic $R$-module $D/A=cR/A$. But by our choice, $B$ is not finitely generated, leading to a contradiction. Thus $P\nsupseteq \Soc(D)=A\oplus B$. Given $A\subseteq P$ and $A\oplus B\nsubseteq P$, there exists a simple submodule $S\subseteq B$ such that $S \nsubseteq P$. Thus $S\cap P=0$ and it implies $\left((S\oplus A)/A\right)\cap \Soc(P/A)=0$. We have
 \[S\cong(S\oplus A)/ A \subseteq \Soc(M)=\Soc(P/A)\oplus \Soc(N/A).\]
 The {semisimple submodule  $\left((S\oplus A)/A\right) \oplus \Soc(P/A) $ of $ M $ must be a direct summand of $ \Soc(M) $. It follows that} $\Soc(M)= ((S\oplus A)/A) \oplus  \Soc(P/A) \oplus (U/A)$ for some submodule $U/A$ of $\Soc(M)$, where $A\subseteq U\subseteq D$. Then $(S\oplus A)/A$ is isomorphic to a direct summand of $\Soc(M)/\Soc(P/A)\cong \Soc(N/A)$. Therefore $N/A$ should have a simple submodule $T/A\cong (S\oplus A)/ A \cong S$. But $S$ is a simple $R$-submodule of $R_R$. Since $R$ is a right V-ring, $S$ is injective. Then it is a direct summand of $R_R$ and hence is projective. Thus $N/A$ has an injective and projective simple submodule $T/A \cong S$ which is a direct summand of $N/A$ (since $S$ is injective). This contradicts the assumption that $N/A$ {is stable}. Therefore, the cyclic $R$-module $M$ does not have a decomposition as a direct sum of a projective submodule and a stable submodule.
\end{proof}

\end{document}
